\documentclass[10pt,a4paper]{amsart}
\usepackage[margin=19mm]{geometry}
\usepackage{amsmath,amssymb,mathtools}
\usepackage[scr=rsfs,cal=euler]{mathalfa}
\usepackage{xfrac}
\usepackage[unicode,hidelinks,bookmarks=false]{hyperref}
\newcommand{\ie}{i.e.\ }
\newcommand{\cf}{cf.\ }
\newcommand{\resp}{respectively\ }
\newcommand{\Subsection}{\subsection}
\providecommand{\nobreakdash}{\nobreak\hbox{-}}
\setlength{\emergencystretch}{2em}
\multlinegap0pt
\usepackage{amssymb}
\usepackage{mathtools}
\usepackage{bm}
\newtheorem{definition}{Definition}
\newtheorem{proposition}{Proposition}
\newcommand{\Ftwo}{\mathbb{F}_2}
\newcommand{\Htil}{\tilde{H}}
\newcommand{\Ker}{\mathrm{Ker}}
\newcommand{\Row}{\mathrm{Row}}
\newcommand{\vect}[1]{\bm{#1}}
\newcommand{\vlambda}{\vect{\lambda}}
\newcommand{\vvct}{\vect{v}}
\newcommand{\vx}{\vect{x}}
\newcommand{\vz}{\vect{z}}
\newcommand{\wt}{\mathrm{wt}}
\title{Heuristic Search for Minimum-Distance Upper-Bound Witnesses in Quantum APM-LDPC Codes}
\author{Kenta Kasai}
\date{Source: arXiv:2604.15307v1 (CC BY 4.0)}
\begin{document}
\maketitle

\noindent\emph{Source and licence.} Heuristic Search for Minimum-Distance Upper-Bound Witnesses in Quantum APM-LDPC Codes. Version arXiv:2604.15307v1 (CC BY 4.0), distributed by arXiv under the Creative Commons Attribution 4.0 International licence, http://creativecommons.org/licenses/by/4.0/, as recorded in the arXiv licence metadata for this exact version and shown on the abstract page https://arxiv.org/abs/2604.15307v1. Full licence notice: \texttt{licenses/arxiv-2604.15307-CC-BY-4.0.txt}.\par\smallskip

\noindent\textit{Editorial scope.} Counted unit: the article's proposition prop:latent-upper-en (source0.tex lines 421--454). The article's complete original proof of it is delivered with it (source0.tex lines 456--503, one \texttt{proof} environment of 48 lines). No mathematics is written, repaired or re-derived: the statement and the proof are the article's own lines, in the article's own notation, and no retained line is substituted or re-flowed.  Retained uncounted as the source-local context the proof needs: lines 226--262.  Nothing was omitted from any retained span, so the package carries no omission record.  No retained line contains a citation, so the wrapper carries no bibliography and no bibliography entry.  No retained line contains a figure environment, an image-inclusion command or drawing code, so no image is delivered and the package declares no asset dependency.  Editorial formatting only, in addition: an AMS \texttt{amsart} preamble that restates the article's own declarations and macros used by the retained lines, namely the environments \texttt{definition}, \texttt{proposition} on separate counters, together with the standard packages those lines need.  The retained environments print in this document as Definition 1, Proposition 1 in order of appearance, whereas the article prints the same items with the numbering of its own full document.  The complete native source of the article as distributed in the arXiv e-print for this exact version is bundled unchanged as \texttt{sources/198599/source0.tex}.
\begin{definition}[Overall, latent, and non-latent distances]
\label{def:latent-nonlatent-distances-en}
From this point on, assume that $H_XH_Z^{\mathsf T}=0$ and define the CSS constituent codes by
\[
C_X:=\Ker(H_X),\qquad C_Z:=\Ker(H_Z).
\]
The usual CSS distances are
\[
d_X=\min\{\wt(\vx):\vx\in C_Z\setminus C_X^\perp\},\qquad
d_Z=\min\{\wt(\vz):\vz\in C_X\setminus C_Z^\perp\}.
\]
To isolate the contribution coming directly from the latent row spaces, define
\[
d_X^{(\mathrm{lat})}:=
\min\{\wt(\vx):\vx\in (C_Z\cap \Row(\Htil_X))\setminus C_X^\perp\},
\]
\[
d_Z^{(\mathrm{lat})}:=
\min\{\wt(\vz):\vz\in (C_X\cap \Row(\Htil_Z))\setminus C_Z^\perp\}.
\]
To measure the complementary contribution outside the latent row spaces, define
\[
d_X^{(\mathrm{nlat})}:=
\min\{\wt(\vx):\vx\in C_Z\setminus C_X^\perp,\ \vx\notin \Row(\Htil_X)\},
\]
\[
d_Z^{(\mathrm{nlat})}:=
\min\{\wt(\vz):\vz\in C_X\setminus C_Z^\perp,\ \vz\notin \Row(\Htil_Z)\}.
\]
All four minima are interpreted as $+\infty$ if the corresponding set is empty.
Then
\[
d_X=\min\{d_X^{(\mathrm{lat})},d_X^{(\mathrm{nlat})}\},
\qquad
d_Z=\min\{d_Z^{(\mathrm{lat})},d_Z^{(\mathrm{nlat})}\}.
\]
\end{definition}

\begin{proposition}[Latent upper bounds from mixed-product kernels]
\label{prop:latent-upper-en}
The latent distances admit the exact parameterizations
\[
d_X^{(\mathrm{lat})}
=
\min\left\{
\wt(\vlambda^{\mathsf T}\Htil_X):
\vlambda\in(\Ftwo^P)^{L/2-J},\ 
H_Z\Htil_X^{\mathsf T}\vlambda=0,\ 
\vlambda^{\mathsf T}\Htil_X\notin C_X^\perp
\right\},
\]
\[
d_Z^{(\mathrm{lat})}
=
\min\left\{
\wt(\vvct^{\mathsf T}\Htil_Z):
\vvct\in(\Ftwo^P)^{L/2-J},\ 
H_X\Htil_Z^{\mathsf T}\vvct=0,\ 
\vvct^{\mathsf T}\Htil_Z\notin C_Z^\perp
\right\},
\]
with the convention that the minimum is $+\infty$ when the corresponding set is empty.

In particular, every feasible coefficient vector
$\vlambda\in(\Ftwo^P)^{L/2-J}$ satisfying
$H_Z\Htil_X^{\mathsf T}\vlambda=0$ and
$\vlambda^{\mathsf T}\Htil_X\notin C_X^\perp$
produces
$\vx:=\vlambda^{\mathsf T}\Htil_X\in (C_Z\cap \Row(\Htil_X))\setminus C_X^\perp$,
hence $d_X^{(\mathrm{lat})}\le \wt(\vx)$ and therefore $d_X\le \wt(\vx)$.
The $Z$-side statement is identical.
\end{proposition}

\begin{proof}
We prove the $X$-side identity. By Definition~\ref{def:latent-nonlatent-distances-en},
\[
d_X^{(\mathrm{lat})}
=
\min\{\wt(\vx):\vx\in (C_Z\cap \Row(\Htil_X))\setminus C_X^\perp\}.
\]
Let
\[
\mathcal{L}_X
:=
\left\{
\vlambda\in(\Ftwo^P)^{L/2-J}:
H_Z\Htil_X^{\mathsf T}\vlambda=0,\ 
\vlambda^{\mathsf T}\Htil_X\notin C_X^\perp
\right\}.
\]

If $\vlambda\in\mathcal{L}_X$ and $\vx:=\vlambda^{\mathsf T}\Htil_X$, then
$\vx\in\Row(\Htil_X)$ and
$H_Z\vx^{\mathsf T}=H_Z\Htil_X^{\mathsf T}\vlambda=0$,
so $\vx\in C_Z$; by definition of $\mathcal{L}_X$, we also have
$\vx\notin C_X^\perp$.
Hence every feasible coefficient vector produces an element of
$(C_Z\cap \Row(\Htil_X))\setminus C_X^\perp$.

Conversely, if $\vx\in (C_Z\cap \Row(\Htil_X))\setminus C_X^\perp$,
then $\vx\in\Row(\Htil_X)$, so there exists some
$\vlambda\in(\Ftwo^P)^{L/2-J}$ with
$\vx=\vlambda^{\mathsf T}\Htil_X$.
Because $\vx\in C_Z$,
$H_Z\Htil_X^{\mathsf T}\vlambda=H_Z\vx^{\mathsf T}=0$,
and since $\vx\notin C_X^\perp$, this $\vlambda$ belongs to $\mathcal{L}_X$.

Therefore the image set
\[
\{\vlambda^{\mathsf T}\Htil_X:\vlambda\in\mathcal{L}_X\}
\]
coincides exactly with
\[
(C_Z\cap \Row(\Htil_X))\setminus C_X^\perp.
\]
This proves the displayed formula for $d_X^{(\mathrm{lat})}$.
Finally, since
$d_X=\min\{d_X^{(\mathrm{lat})},d_X^{(\mathrm{nlat})}\}$,
every latent witness also yields the overall upper bound $d_X\le \wt(\vx)$.
The $Z$-side argument is identical.
\end{proof}

\end{document}
